In the last decade, commercial resistive random-access memory (ReRAM)
modules have entered mass production as embedded non-volatile data
storage for low-power microcontrollers and edge devices. ReRAM write
operations complete through stochastic program-verify loops. The
duration of each loop depends on the physical state of individual memory
cells, making ReRAM an attractive entropy source for security primitives
such as true random number generators and physically unclonable
functions. Commercial devices improve reliability by augmenting user
data with error-correcting codes that introduce additional switching
activity on every write. Together, these mechanisms create structured
timing leakage channels that have remained uninterpretable without
costly, destructive reverse engineering of the underlying architecture.
As a result, timing models are inaccurate, the security primitives that
rely on ReRAM switching entropy cannot be properly characterized, and
the timing side-channel attack surface of commercial ReRAM products
remains uncharacterized.

In crossbar architectures, strict sequencing constraints on SET and
RESET operations produce data-dependent write latency distributions.
Prior investigations of commercial ReRAM timing identified structured
variance in these distributions but could not explain or model it.
Embedded, undocumented error-correcting codes (ECCs) are the dominant
source of this structured variance. The timing fluctuations introduced
by ECC switching activity is large enough to obscure the stochastic
device physics that prior security and entropy analyses assumed they
were measuring. This dissertation closes that gap through systematic
experimental characterization of commercial ReRAM write latency,
revealing these dominant structural factors and demonstrating that they
are recoverable through non-invasive side-channel analysis.

We introduce a novel technique for \underline{B}it-exact
\underline{R}ecovery of \underline{E}CC from \underline{W}rite Timing in
\underline{R}eRAM \underline{C}rossbars (BREW-RC), the first technique
to recover the bit-exact ECC parity submatrix from commercial ReRAM
memories using only write-access timing, without physical probing or
device modification. BREW-RC formulates ECC recovery as a SAT problem,
deriving constraints from observed write latency distributions and
recovering the complete generator matrix from two commercial products.
The recovered matrix enables accurate codeword-level timing models,
substantially improving side-channel analysis and enabling non-invasive
reliability characterization of devices with undisclosed error
correction capabilities.

We apply this structural understanding to characterize stochastic write
latency. Prior work on commercial ReRAM entropy sources applied
aggressive debiasing to suppress structured variance of unknown origin.
We demonstrate that this structure is fully explained by ECC activity.
Isolating the deterministic ECC contribution from the underlying
stochastic filament dynamics enables stationary, wear-leveled
measurement sequences with min-entropy rigorously estimated under NIST
SP 800-90B. We then characterize how device aging affects switching
dynamics and entropy yield. The resulting TRNG passes NIST SP 800-22 and
PractRand and achieves a 17x throughput improvement over prior ECC- and
wearout-unaware methods.

% Finally we describe how these contributions connect to an emerging
% class of hardware security threats in which the physical state of
% deployed devices encodes exploitable information that their logical
% interfaces do not expose. FPGA Pentimenti, which recover prior device
% data through Bias Temperature Instability effects on programmable
% routing, are the most developed instance of this class. The timing
% channels characterized in this dissertation occupy an analogous
% position for ReRAM, establishing the observability, structural
% modeling, characterization that exploitable data remanence in
% commercial ReRAM will require.

% % %% %% %% %% %% %% %% %% %% %% %% %% %% %% %% %% %% %% %% %% %% %% %%

% % Prior Pentimenti-heavy Framing 

% %% %% %% %% %% %% %% %% %% %% %% % %% %% %% %% %% %% %% %% %% %% %% %%

% Abstract Recent discoveries have demonstrated a class of security
% vulnerabilities that violate isolation guarantees in FPGA cloud systems
% through leakage channels emerging from low-level transistor device
% physics, known as FPGA Pentimenti. Historically, Pentimenti are traces
% of an earlier composition visible beneath the surface of a painting. In
% the context of hardware security, FPGA Pentimenti are remnants of prior
% device state that persist after logical erasure, manifesting as
% threshold voltage shifts driven by Bias Temperature Instability (BTI)
% at transistor gate oxide interfaces. Pentimento-based attacks have been
% demonstrated on Xilinx FPGAs and require high-precision propagation
% delay sensors, the ability to propagate measurement pulses through
% sensitive assets within the FPGA fabric, and the ability to integrate
% BTI effects across many of these assets (i.e. by chaining routing
% resources). These requirements are practical on FPGAs where highly
% reconfigurable fabrics allow attackers to freely place propagation
% delay sensors and perform measurements on vulnerable assets. However,
% it remains an open research challenge to find Pentimenti-like threats
% outside of FPGA routing resources.

% Resistive random access memories (ReRAMs) expose leakage channels that
% resemble those exploited to recover FPGA Pentimenti. ReRAM devices
% leverage non-deterministic programming sequences which operate on
% stochastic device physics. Measuring the duration of these sequences
% enables the construction of leakage channels which may expose the
% internal state of ReRAM devices, proprietary architectural details,
% and, perhaps, remnants of prior states (i.e. ReRAM Pentimenti). In the
% last decade, commercial ReRAM products have entered mass production
% making it possible to investigate timing side-channels that may expose
% these leakage channels. However, commercial ReRAM timing is effected
% significantly by the underlying crossbar memory architecture, the added
% switching activity of error correcting codes, and the wearout state of
% individual ReRAM cells. These factors introduce significant gaps that
% must be solved before the deployment of Pentimenti-like attacks in
% commercial ReRAM are possible.

% This work presents the leakage channels necessary for the study of
% Commercial ReRAM Pentimenti. First, we characterize the architectural
% timing leakage channels in commercial ReRAM crossbar memories,
% identifying how polarity serialization in the crossbar write path
% produces data-dependent write latency distributions and demonstrating
% that embedded, undocumented error correcting codes (ECCs) are the
% dominant source of structured timing variance in commercial products.
% We introduce BREW-RC, the first technique to recover the bit-exact ECC
% parity submatrix from commercial ReRAM memories using only write-access
% timing, without physical probing or device modification. Second, we
% demonstrate that knowledge of the recovered ECC is necessary but not
% sufficient to characterize the residual stochastic component of ReRAM
% write latency. By applying ECC-aware analysis to isolate deterministic
% parity-driven timing perturbations from the underlying stochastic
% filament dynamics, we construct a framework for evaluating the true
% entropy of commercial ReRAM timing sources, validated against NIST SP
% 800-90B and shown to improve entropy harvesting throughput by 17x over
% prior ECC-unaware methods. Together, these contributions establish the
% observability and structural understanding required to pursue
% Commercial ReRAM Pentimenti: identifying the timing leakage channels,
% characterizing their deterministic and stochastic components, and
% isolating the device-state remnants that may encode the analog
% remanence of prior ReRAM cell states.
